\subsection*{Causal and Normative Concepts}
\textit{Terms relating to the classes of commitment excluded from neutral substrates.}

\begin{description}[style=nextline,leftmargin=1.5cm]
  \item[Causal Attribution]
        The assignment of a causal relation between events,
        entities, or occurrences within an interpretive framework.
        Causal attribution is typically framework-dependent:
        since different causal models may attribute causal influence differently.
        Contrast with \emph{Causal Commitment}.

  \item[Causal Commitment]
        An ontological commitment that asserts a causal relation
        (e.g., $A$ caused $B$) as a substrate-layer fact.
        Causal commitments depend on the causal model or explanatory framework
        under which the relation is interpreted and may therefore vary
        across admissible causal frameworks.
        In a neutral substrate, such commitments cannot be asserted
        at the foundational layer.

  \item[Deontic Commitment]
        An ontological commitment to obligations, permissions, prohibitions,
        or duties as part of the ontology's representation.
        Deontic commitments are a subset of normative commitments
        and encode what ought to be done rather than what exists or occurs.
        They may vary across admissible normative frameworks.
        In a neutral substrate, such commitments cannot be asserted
        at the foundational layer.

  \item[Deontic Predicate]
        A predicate symbol whose extension is defined by
        obligations, permissions, or prohibitions.
        Deontic predicates are normative in character
        and are excluded from the foundational Level of Abstraction
        of a neutral substrate.

  \item[Framework-Dependent]
        A property of claims or predicates whose truth value
        may vary across admissible interpretive frameworks.
        Causal and normative claims are typically framework-dependent.
        In the domains considered in this paper, identity and persistence
        conditions are assumed to remain invariant across admissible frameworks;
        if they are contested, neutral substrates are not attainable
        under the neutrality definition.

  \item[Normative Attribution]
        The assignment of obligations, violations, permissions,
        or responsibilities to an entity or occurrence
        by an interpretive framework.
        Normative attribution is framework-dependent
        and therefore appears in interpretive layers rather than
        as substrate-layer commitments.
        Contrast with \emph{Normative Commitment}.

  \item[Normative Commitment]
        An ontological commitment that encodes evaluative, legal, moral,
        or policy-based judgments as part of the foundational layer.
        Normative commitments may reflect ethical, legal, institutional, or procedural
        frameworks and therefore may vary across admissible interpretive perspectives.

  \item[Pre-Causal]
        The property of an ontology whose foundational layer excludes
        causal commitments.
        In a neutral substrate, this also requires that the foundational
        Level of Abstraction exclude causal predicates.
        Pre-causal does not mean acausal or causation-denying;
        it means \emph{prior to causal attribution}.
        A pre-causal substrate represents entities and occurrences
        without asserting causal relations as substrate-layer commitments,
        allowing causal reasoning to occur in interpretive layers
        where model assumptions are explicit and disagreement is representable.
        Causation is externalized, not eliminated.

  \item[Pre-Normative]
        The property of an ontology whose foundational layer excludes
        normative commitments.
        In a neutral substrate, this also requires that the foundational
        Level of Abstraction exclude normative predicates.
        Pre-normative does not mean non-normative or norm-denying;
        it means \emph{prior to normative attribution}.
        A pre-normative ontology represents entities and occurrences
        without asserting obligations, permissions, or evaluations
        as substrate-layer commitments,
        allowing normative reasoning to occur in interpretive layers
        where framework assumptions are explicit and disagreement is representable.

\end{description}
